% Part: normal-modal-logic
% Chapter: syntax-and-semantics
% Section: normal-modal-logics

\documentclass[../../../include/open-logic-section]{subfiles}

\begin{document}

\olfileid[hi]{nml}{syn}{nor}

\olsection{सामान्य मोडल तर्कशास्त्र}

किसी मोडल तर्कशास्त्र को !!{formula}ों का कोई उपयुक्त रूप से सुव्यवस्थित
समुच्चय माना जा सकता है। कुछ मोडल तर्कशास्त्र विशेष रुचि के हैं---इसलिए कि
उनकी रोचक व्याख्याएँ और अनुप्रयोग हैं, पर मुख्यतः इसलिए कि उन्हें सैद्धांतिक
रूप से अच्छी तरह समझा गया है। ये तथाकथित सामान्य मोडल तर्कशास्त्र हैं।

\begin{defn}\ollabel{defn:modal-logic}
\emph{मोडल तर्कशास्त्र}~$\Log L$, मूल मोडल भाषा के !!{formula}ों का ऐसा
समुच्चय है जो
\begin{enumerate}
\item सभी प्रतिज्ञप्तिक सर्वसत्यताओं को समाहित करता है,
\item एकरूप प्रतिस्थापन के अधीन संवृत है, और
\item मोडस पोनेन्स के अधीन संवृत है; अर्थात् जब भी $!A$ और
  $(!A \lif !B) \in \Log L$, तब $!B \in \Log L$ भी।
\end{enumerate}
\end{defn}

\begin{ex}
!!{formula}ों के कुछ ऐसे समुच्चय, जो निस्संदेह बहुत रोचक नहीं हैं, पर इस
परिभाषा को संतुष्ट करते हैं और इसलिए मोडल तर्कशास्त्र माने जाते हैं:
\begin{enumerate}
\item सभी सर्वसत्यताओं का समुच्चय;
\item सभी !!{formula}ों का समुच्चय (असंगत मोडल तर्कशास्त्र)।
\end{enumerate}
\end{ex}

\begin{defn}\label{defn:normal-modal-logic}
कोई मोडल तर्कशास्त्र~$\Log L$ \emph{सामान्य} है तभी और केवल तभी जब
\begin{enumerate}
\item उसमें ये दोनों हों:
\begin{align}
\tag{K} \Box(p \lif q) & \lif (\Box p \lif \Box q) \\
\tag{Dual} \Diamond p & \liff \lnot \Box \lnot p;
\end{align}
\item और वह आवश्यकीकरण के अधीन संवृत हो; अर्थात् जब भी
  $!A \in \Log L$, तब $\Box !A \in \Log L$ भी।
\end{enumerate}
\end{defn}

हमें मोडल तर्कशास्त्र परिभाषित करने के मुख्यतः दो ढंग मिलेंगे। एक
प्रमाण-सैद्धांतिक ढंग है: कुछ !!{derivation}-तंत्रों में व्युत्पन्न किए जा
सकने वाले !!{formula}ों का समुच्चय। दूसरा प्रतिरूप-सैद्धांतिक ढंग है: प्रतिरूपों
के कुछ वर्गों में मान्य !!{formula}ों का समुच्चय। यह साधारण प्रतिज्ञप्तिक
तर्कशास्त्र (और इसी तरह प्रथम-कोटि तर्कशास्त्र) की स्थिति को प्रतिबिंबित करता
है। वहाँ भी !!{formula}ों के किसी समुच्चय को दो भिन्न ढंगों से चुना जा सकता
है, जैसे सभी सत्य-मूल्य नियतनों में सत्य सर्वसत्यताओं का समुच्चय, और किसी
!!{derivation}-तंत्र में सभी !!{derivable} !!{formula}ों का समुच्चय। सामान्य
मोडल तर्कशास्त्रों के लिए मोडल तर्कशास्त्र निर्दिष्ट करने का यह द्विविध ढंग
संबंधात्मक प्रतिरूपों और अभिगृहीतीय (हिल्बर्ट-प्रकार) प्रमाण-तंत्रों से संभव
है।
\end{document}
